\documentclass[11pt]{article}
\usepackage[margin=1.1in]{geometry}
\usepackage{amsmath, amssymb, booktabs}
\usepackage[colorlinks=true, linkcolor=blue]{hyperref}

\newcommand{\E}{\mathbb{E}}
\newcommand{\Post}{\mathrm{Post}}

\title{The Role of Institutions in the Output Response:\\ What We Do and What We Find}
\author{}
\date{\today}

\begin{document}
\maketitle

\noindent
Companion to \texttt{inst\_te\_specification.tex}, which derives the estimand
and the identification argument. This document reports the empirical pipeline
and results. All numbers come from \texttt{inst\_fe\_diag.do}.

\section{The exercise in one paragraph}

We ask whether the institution a PI sits at shaped how her output responded
to the 2014 procurement shock, and if so, which institutional
characteristics account for it. Because no PI in the sample changes
institution, an institution \emph{level} effect is collinear with the PI
fixed effect and is not identified. What is identified is the institution
$\times$ \Post{} \emph{shift}: two PIs with the same exposure, one at
institution $j$ and one at $k$, differ in their post-2014 log-output change
by exactly $te_j - te_k$. Stage 1 estimates the full set of $te_j$ from one
pooled PPML; Stage 2 treats the $\widehat{te}_j$ as noisy observations of
true institution effects and asks, by Fay--Herriot meta-regression, which
institutional characteristics predict them.

\paragraph{Stage 1.}
\[
\E[Y_{it}] = \exp\Big( \alpha_i + \delta_t + b\,Z_{it} + g\,Z^{S}_{it}
+ \textstyle\sum_{j \neq j_0} te_j\,\mathbf{1}[j(i)=j]\cdot \Post_t \Big),
\]
by \texttt{ppmlhdfe}, PI and year absorbed, standard errors clustered by PI.
The exposure response $b$ is held common across institutions, so $te_j$ is a
shift \emph{conditional on the same exposure}. The reference institution
$j_0$ is the largest (Johns Hopkins, 464 PIs); estimates are then re-centered
on the output-weighted mean, so zero is the response of the average unit of
research output.

\paragraph{Stage 2.}
The $\widehat{te}_j$ are individually noisy, so they are never used raw as
data points. They enter a random-effects meta-regression that models their
sampling error explicitly, weighting institution $j$ by
$1/(\tau^2 + \widehat{se}_j^2)$, with $\tau^2$ the variance of the true
institution effects estimated by REML and Knapp--Hartung standard errors.
Empirical Bayes posterior means are computed for display only; no reported
slope or variance decomposition uses them.

\section{Sample}

\begin{center}
\begin{tabular}{@{}lrl@{}}
\toprule
Step & $N$ & Binding constraint \\
\midrule
PI-by-year panel, 2010--2019 & 123,370 & 12,337 PIs, 198 institutions \\
PIs changing institution     & 0        & sample construction \\
Institutions with an estimable $te_j$ & 196 & reference + 1 collinear \\
\quad and $\ge 10$ PIs on the estimation sample & 128 & precision floor \\
\quad and endowment + HERD characteristics      & 103 & IPEDS/HERD coverage \\
\quad and all 24 characteristics                & 98  & complete cases \\
\bottomrule
\end{tabular}
\end{center}

\noindent
The attrition from 198 to 103 is coverage, not selection on outcomes, but it
does skew the second stage toward larger institutions. The pooled exposure
response in the Stage 1 fit is $\hat b = -0.249$ (se $0.130$), consistent
with the main reduced form.

\section{Finding 1: institutions matter, but modestly}

A joint Wald test that every institution $\times$ \Post{} shift is zero gives
$\chi^2(187) = 14{,}841$, decisively rejected. Homogeneity across the 128
estimable institutions is also rejected, $Q = 161.5$ on 127 degrees of
freedom ($p = 0.021$).

The magnitude, however, is modest. On the horse-race sample the REML estimate
is $\tau^2 = 0.0024$, so the standard deviation of \emph{true} institution
effects is $\tau \approx 0.049$ log points, and $I^2 = 27.8\%$. Moving from
an average institution to one a standard deviation better is worth about five
log points of post-2014 output --- real, but an order of magnitude smaller
than the pooled exposure response.

There is a size gradient visible in the normalization itself. The unweighted
mean $te_j$ is $-0.111$ while the output-weighted mean is $-0.074$: the
average \emph{unit of science} fared better than the average
\emph{institution}, because the institutions where most output happens are
the large ones, and large institutions did better.

\section{Finding 2: most of the raw spread is estimation error}

The raw $\widehat{te}_j$ have standard deviation $0.116$ across the 128
institutions, but the median sampling standard error is $0.092$. Decomposing,
$0.116^2 \approx 0.044^2 + 0.107^2$: roughly \textbf{85\% of the observed
dispersion is noise}, and the true dispersion is the $\tau$ of Finding 1.
This is why the empirical Bayes shrinkage factor
$\lambda_j = \tau^2/(\tau^2 + \widehat{se}_j^2)$ averages only $0.21$,
never exceeding $0.47$ even for the best-measured institution.

Two practical consequences. First, no individual institution's $te_j$ should
be quoted or ranked: an estimate at $-0.40$ is far more likely to be an
ordinary institution with an unlucky draw than a genuinely extreme one.
Second, neither caterpillar plot shows ``how much institutions differ.'' The
raw figure overstates it by a factor near three; the EB figure understates
it, because posterior means are under-dispersed by construction
($\mathrm{sd}(\widehat{te}^{EB}) = \sqrt{\bar\lambda}\,\tau = 0.020$). Only
$\tau = 0.049$ answers that question.

\section{Finding 3: the characteristic that matters is the scale of biomedical research}

Every one of the 31 characteristics in the univariate screen has a
\emph{positive} dose-adjusted coefficient: larger and better-funded
institutions had less negative post-2014 responses. The ranking, per standard
deviation of the characteristic, with Benjamini--Hochberg $q$:

\begin{center}
\begin{tabular}{@{}lrr@{}}
\toprule
Characteristic (log) & $\hat\theta$ & $q_{BH}$ \\
\midrule
Non-federal life-science/biomedical funding & $+0.058$ & $2\times10^{-6}$ \\
Total biomedical funding                    & $+0.056$ & $3\times10^{-6}$ \\
Federal life-science/biomedical funding     & $+0.051$ & $2\times10^{-5}$ \\
Total R\&D funding                          & $+0.049$ & $5\times10^{-5}$ \\
Basic research expenditure                  & $+0.044$ & $6\times10^{-5}$ \\
Endowment per FTE                           & $+0.031$ & $4\times10^{-4}$ \\
\midrule
\multicolumn{3}{@{}l@{}}{\emph{not distinguishable from zero:}} \\
Funding diversity ($1-$HHI over HERD sources) & $+0.008$ & $p = 0.35$ \\
Development expenditure                       & $+0.004$ & $p = 0.65$ \\
Health-services funding                       & $+0.013$ & $p = 0.24$ \\
\bottomrule
\end{tabular}
\end{center}

\noindent
The top rows survive Bonferroni, not merely BH. The nulls are informative:
this is not a story about diversified revenue cushioning the shock, nor about
applied-versus-basic orientation. It is about the \emph{level} of basic
biomedical research resources.

In the three-characteristic horse race, endowment and total R\&D funding both
survive jointly while life-science funding drops out as the collinear
duplicate:

\begin{center}
\begin{tabular}{@{}lrrr@{}}
\toprule
 & $\hat\theta$ & se & $p$ \\
\midrule
Endowment per FTE (z)      & $+0.0217$ & $0.0090$ & $0.018$ \\
Total R\&D funding (z)     & $+0.0530$ & $0.0201$ & $0.010$ \\
Life-science funding (z)   & $-0.0158$ & $0.0197$ & $0.426$ \\
Institution mean exposure (z) & $-0.0270$ & $0.0118$ & $0.024$ \\
\bottomrule
\end{tabular}
\end{center}

\noindent
Because the funding measures are near-collinear, individual $\hat\theta$
should be read ordinally. The collinearity-proof version is the principal
components companion: PC1 carries 71\% of the variance of the 24
characteristics, loads positively on every funding measure --- it is the scale
factor --- and enters at $+0.0134$ ($p < 0.001$). The 24-characteristic joint
horse race produces no individually significant coefficient, exactly as
collinearity predicts, while still driving residual heterogeneity to zero.

\section{Finding 4: exposure concentration hurts, beyond a PI's own exposure}

The institution's mean PI exposure enters with a robustly negative
coefficient: $-0.026$ alone ($p = 0.041$), $-0.027$ in the horse race
($p = 0.024$), $-0.028$ once baseline productivity is added ($p = 0.017$).
Since each PI's \emph{own} exposure is already absorbed by the pooled $b
Z_{it}$, this says that conditional on how exposed you were, being at an
institution where your colleagues were also heavily exposed made your
response worse. The leave-one-out version in the one-step check --- which
excludes the PI's own exposure from the institution mean by construction ---
confirms it at $-0.019$ ($p = 0.011$).

This is the one result that is not a mechanical size story, and it is the
natural place to argue for a within-institution spillover: shared equipment,
shared core facilities, or shared administrative capacity absorbing a common
procurement shock.

\section{Finding 5: observables explain essentially all of the heterogeneity}

\begin{center}
\begin{tabular}{@{}lrr@{}}
\toprule
Specification & $\tau^2$ & $Q_{res}$ $p$ \\
\midrule
(a) no covariates          & $0.00242$ & $0.008$ \\
(b) dose only              & $0.00196$ & $0.018$ \\
(c) characteristics + dose & $\approx 0$ & $0.382$ \\
\bottomrule
\end{tabular}
\end{center}

\noindent
Scale plus exposure concentration absorb the entire estimated between-institution
variance, and the residual homogeneity test stops rejecting. There is no
detectable residual ``institution quality'' once these are conditioned on.
Adding pre-2014 institution output (Finding 3's characteristics compete with
it: $+0.022$, $p = 0.051$) leaves the dose coefficient untouched but
attenuates endowment and funding, as expected for a third size proxy.

\section{Validation}

The two-step procedure is cross-checked against a one-step PPML in which
$\Post_t \times W_k$ replaces the estimated $te_j$ entirely, with the same
fixed-effect structure. The two approaches agree closely:

\begin{center}
\begin{tabular}{@{}lrr@{}}
\toprule
 & Two-step $\hat\theta_k$ & One-step $\hat\delta_k$ \\
\midrule
Endowment per FTE    & $+0.0217$ & $+0.0193$ \\
Total R\&D funding   & $+0.0530$ & $+0.0564$ \\
Life-science funding & $-0.0158$ & $-0.0276$ \\
Dose                 & $-0.0270$ & $-0.0194$ \\
\bottomrule
\end{tabular}
\end{center}

\noindent
Same signs, same magnitudes, significant in the same places. The findings are
not an artifact of estimating institution effects first and regressing second.

\section{Caveats}

\begin{enumerate}
\item \textbf{Collinearity.} All funding measures move together; the honest
claim is ``scale of biomedical research,'' not any specific funding stream.
\item \textbf{Sorting.} With no movers, $te_j$ cannot separate a causal
institution effect from the resilience of the PIs an institution attracts.
The effect is descriptive.
\item \textbf{Coverage.} The second stage runs on 103 of 198 institutions,
skewed large.
\item \textbf{$\tau^2$ is probably understated.} The $\widehat{se}_j$ are
measured relative to the omitted reference, so they carry that institution's
estimation error, which is common to all $te_j$ and contributes nothing to
cross-institution dispersion. REML subtracts it anyway. This biases $\tau^2$
down, shrinkage up, and makes the ``100\% explained'' in Finding 5 look
cleaner than it is. The Wald statistic of 14,841 against an expected $\approx
320$ points the same way. The fix is to feed REML the diagonal of $MVM'$ with
$M = I - \mathbf{1}\mathbf{1}'/J$; the independent check is a split-sample
estimate of $\tau^2$.
\item \textbf{Screening-grade inference.} Second-stage standard errors are
Knapp--Hartung on PI-clustered $\widehat{se}_j$; anything headline should be
bootstrapped through both stages.
\item \textbf{PC5.} Significant at $p<0.001$ while carrying 2.5\% of the
variance --- treat as overfitting until shown otherwise.
\end{enumerate}

\section{Bottom line}

Institutional heterogeneity in the post-shock output response is
statistically unambiguous but quantitatively modest: the true standard
deviation across institutions is about five log points, against a pooled
exposure response of $-25$. That heterogeneity is fully accounted for by two
things --- the scale of an institution's biomedical research enterprise, which
helped, and the concentration of exposure among its PIs, which hurt --- with
no residual institution-specific component. Institutions mattered, but
through resources and shared exposure rather than through anything
idiosyncratic to the place.

\end{document}
